% =========================================================
% Rational Cauchy Sequences
% =========================================================

\subsection{Rational Cauchy Sequences}
\label{sec:rational-cauchy}

\subsubsection{Definitions and Theorems}

The rational numbers form an ordered field and they approximate real numbers densely.  At this
stage it is tempting to believe that they should already be adequate for analysis.  The first serious
obstruction appears when one studies limits of rational sequences.

\begin{definition}[Sequence of rational numbers]
\label{def:rational-sequence}
A \emph{sequence of rational numbers} is a function
\[
x:\mathbb{N}\to\mathbb{Q}.
\]
It is usually written $(x_n)_{n\in\mathbb{N}}$ or simply $x_1,x_2,x_3,\dots$.
\end{definition}

\begin{remark}
Thinking of a sequence as a function is useful because it emphasizes that a sequence is not merely a
list but an indexed family of values.  Later this viewpoint makes it natural to add, subtract, and
compare sequences term by term.
\end{remark}

\begin{definition}[Cauchy sequence in $\mathbb{Q}$]
\label{def:cauchy-rational}
A sequence $(x_n)$ of rational numbers is called a \emph{Cauchy sequence} if
\[
\forall \varepsilon>0\ \exists N\in\mathbb{N}\ \forall m,n\ge N,
\qquad
|x_m-x_n|<\varepsilon.
\]
\end{definition}

\begin{remark}
A Cauchy sequence is one whose terms eventually cluster together as tightly as we wish.  Notice that
this definition does not mention a limit at all; it speaks only about the internal behavior of the
sequence.
\end{remark}

\begin{theorem}
\label{thm:convergent-implies-cauchy-q}
Every convergent sequence in $\mathbb{Q}$ is a Cauchy sequence.
\end{theorem}

\begin{proof}
Suppose $x_n\to x\in\mathbb{Q}$.  Let $\varepsilon>0$.  Choose $N$ such that
\[
n\ge N \Rightarrow |x_n-x|<\varepsilon/2.
\]
Then if $m,n\ge N$,
\[
|x_m-x_n|\le |x_m-x|+|x_n-x|<\varepsilon.
\]
Hence $(x_n)$ is Cauchy.
\end{proof}

\begin{remark}
The converse is the issue.  In a complete space every Cauchy sequence converges.  The rational
numbers are not complete, and so this converse can fail.
\end{remark}

\begin{remark}[Examples]
\begin{enumerate}
  \item The sequence $x_n=1/n$ is Cauchy because it converges to $0$.
  \item The decimal approximations
  \[
  1,\ 1.4,\ 1.41,\ 1.414,\ 1.4142,\dots
  \]
  form a Cauchy sequence in $\mathbb{Q}$.
\end{enumerate}
\end{remark}

\begin{theorem}
\label{thm:cauchy-not-convergent-q}
There exist Cauchy sequences in $\mathbb{Q}$ that do not converge to a rational number.
\end{theorem}

\begin{remark}
The second example above converges to $\sqrt{2}$ in $\mathbb{R}$, but $\sqrt{2}\notin\mathbb{Q}$.  Thus
$\mathbb{Q}$ already contains enough arithmetic to define the sequence, but not enough numbers to
contain its limit.
\end{remark}

\subsubsection{Consequences}

\begin{remark}[Motivation for the real numbers]
\label{rem:real-motivation}
One standard way to construct the real numbers is to take equivalence classes of rational Cauchy
sequences, where two Cauchy sequences are considered equivalent if their difference tends to $0$.
In this way the missing limits are formally adjoined to the rationals.
\end{remark}
